% TOP_PROBLEM: {"id": 20001122, "title": "Finite-field point-line incidences and a projective no-wrap criterion", "category_id": 2, "category": {"id": 2, "name": "combinatorics", "display_name": "Combinatorics", "description": "Counting problems, graph theory, discrete structures.", "slug": "combinatorics", "order_index": 2, "created_at": "2026-07-31T15:26:25.671Z"}, "status": "partially_solved", "problem_number": "AIM-COMBINATORICS-0247", "set_id": 14, "statusQualification": "Explicitly isolates the logarithmic-size subquestion; the unrestricted global incidence inequality in the source is false, as its own remarks explain."}
% FIRST_POSTED: 2026-10-01
% ENTRY_KIND: statement_only
% UnsolvedMath ID 20001122; source catalogue code AIM-COMBINATORICS-0247
% !TeX program = lualatex
% Definition and sources only; this file is not an LLM solution attempt.
\documentclass[11pt,a4paper]{article}
\usepackage[margin=25mm]{geometry}
\usepackage{fontspec}
\setmainfont{lmroman10-regular.otf}[BoldFont=lmroman10-bold.otf,
 ItalicFont=lmroman10-italic.otf,BoldItalicFont=lmroman10-bolditalic.otf]
\usepackage{amsmath,amssymb,mathtools}
\usepackage{unicode-math}
\setmathfont{latinmodern-math.otf}
\AtBeginDocument{\let\setminus\smallsetminus}
\usepackage{xurl,hyperref}
\hypersetup{colorlinks=true,urlcolor=blue}
\setcounter{secnumdepth}{0}
\setlength{\parindent}{0pt}
\setlength{\parskip}{5pt}
\setlength{\emergencystretch}{3em}
\begin{document}
\section*{Finite-field point-line incidences and a projective no-wrap criterion}\label{20001122}
{\small UnsolvedMath ID 20001122 · AIM-COMBINATORICS-0247 · Combinatorics · AIM Workshop Problem Lists}
\subsection{Definitions and mathematical statement}
Let $p$ be a prime, let $P\subseteq\mathbb F_p^2$ be a set of distinct points, and let $\mathcal L$ be a set of distinct affine lines.
A line is a set $\{(x,y):ax+by=c\}$ with $(a,b)\ne(0,0)$, with scalar multiples of an equation representing the same line.
Write
\[
 I(P,\mathcal L)=|\{(v,\ell)\in P\times\mathcal L:v\in\ell\}|,
 \qquad n=|P|,\quad l=|\mathcal L|.
\]
The source asks for a prime-field analogue of the real point--line incidence bound and singles out logarithmically small configurations.
Its unrestricted proposed bound has a large-configuration obstruction, so the precise remaining target here is the following.

For every fixed pair of constants $0<c\leq C<\infty$, do there exist $K(c,C)$ and $p_0(c,C)$ such that every prime $p\geq p_0$ and every configuration with
\[
 c\log p\leq n,l\leq C\log p
\]
satisfies
\[
 I(P,\mathcal L)\leq K(c,C)(nl)^{2/3}?
\]
The logarithm is natural, and the constant must be independent of the prime and the point and line choices.
In this size regime the usual additive terms $n+l$ can be absorbed in the displayed estimate.
No general-position assumptions are imposed: many points may be collinear, and many lines may pass through a point.
Only prime fields are considered; extension fields introduce additional subfield configurations.
\subsection{Short English statement}
Does the real-plane incidence exponent remain valid over a prime field when both the point set and line set have logarithmic size?
\subsection{Sources}
\begin{itemize}
\item[S1] ULAM.AI and the UnsolvedMath Contributors, supplied problems.json, record 20001122 (AIM-COMBINATORICS-0247), AIM Workshop Problem Lists. Catalogue identity and source transcription; CC BY 4.0.
\url{https://huggingface.co/datasets/ulamai/UnsolvedMath}.
\item[S2] American Institute of Mathematics, Recent trends in additive combinatorics, item 3.10. Original workshop question underlying AIM-COMBINATORICS-0247.
\url{https://aimath.org/WWN/additivecomb/additivecomb.pdf}.
\end{itemize}
\end{document}
